% !TEX root = ../../tree_notes.tex
% =====================================================================
%  4.3  Encodings and Input Formats                       sec:c04-formats
%  Exports: def:c04-encoding  def:c04-levelorder  thm:c04-selfdelim
%           lst:c04-formats  ex:c04-encodings
%           prob:c04-string-tree    prob:c04-recover-preorder
%           lst:c04-string-tree     lst:c04-recover-preorder
% =====================================================================
\section{Encodings and Input Formats}\label{sec:c04-formats}

\begin{previouslybox}
\begin{enumerate}[nosep, leftmargin=1.6em]
  \item What single loop turns \code{parent[]} into \code{children[]}? (\coderef{c04-reps})
  \item Why does a DFS over an \emph{undirected} adjacency list need a parent guard, and
        what does rooting give you for free? (\thmref{c04-dfs-orient})
  \item In LC 331 (\probref{c03-verify-pre}) what made a preorder string decodable without
        rebuilding -- and how is that the same word ``self-delimiting'' used here?
\end{enumerate}
\end{previouslybox}

\begin{topicintro}
Before there is a tree, there is a \emph{string} -- or an array, or a list of edges. The
judge never hands you pointers; it hands you an \term{encoding}, and your first job is to
\term{decode} it into a structure you can traverse. This section is a field guide to the
encodings contests actually use: LeetCode's \textbf{level-order array with null
marks}, the \textbf{preorder string with sentinels}, the \textbf{nested bracket} (Lisp-like)
form, and the \textbf{depth-dashed preorder} of LC 1028. Each is a different way of laying
the same tree flat; each needs a different reader. You will also \emph{encode} (produce such
a string), which is where LC 606 lives, and see exactly which encodings are
\textbf{uniquely} reversible (\thmref{c04-selfdelim}).
\textbf{Prerequisites.} \secref{c03-dfs} and \secref{c03-rebuild} (preorder/inorder,
serialisation); \secref{c04-reps} (the structures you decode \emph{into}); a stack.

\textbf{After this section you can}
\begin{itemize}[nosep]
  \item read a level-order array (with \# sentinels) into a pointer tree and back;   \hfill\lv{3}
  \item emit the parenthesis encoding of a binary tree under LC 606's rules; \hfill\lv{2}
  \item parse a nested bracket / depth-dashed string with a stack or recursion;      \hfill\lv{3}--\lv{4}
  \item state why a bare preorder is \emph{not} uniquely decodable without marks.     \hfill\lv{4}
\end{itemize}
\end{topicintro}

\begin{recognitionbox}[Pattern card -- the tree arrives as a string/array]
\reckey{the input is \emph{one} $\texttt{string}$ or a $\texttt{vector<string>}$ / array that
\emph{encodes} the tree -- dashes, parentheses, commas with \#, or a preorder with depths --
and you must reconstruct (or produce) the structure before solving.}
\recuse{
  \item ``parse this expression / nested list / serialised tree'' -- the string \emph{is} the
        input; decode it (\probref{c04-recover-preorder}, LC 736);
  \item ``output the tree in this parenthesised form'' -- \emph{encode} it (\probref{c04-string-tree});
  \item LeetCode's \code{deserialize} helpers -- a level-order array with nulls.
}
\recnotuse{
  \item the input is $n{-}1$ edges or a \code{parent[]} array -- that is a \emph{representation},
        \secref{c04-reps}/\secref{c04-rooting}, not a string to parse;
  \item you are given two \emph{traversal arrays} and must rebuild -- LC 105/106, that is
        reconstruction, \secref{c03-rebuild};
  \item the tree is fully built and the string is only the \emph{answer format} -- a formatting
        task, not an encoding.
}
One-line test: \emph{is my input (or output) a flat string/array whose layout, not just
contents, carries the tree shape? If yes, decode/encode -- this section.}
\end{recognitionbox}

\subsection{Intuition}

Here is running example $R$ (\figref{c04-running-r}) written five ways. The tree never
changed -- only the \emph{contract} for reading it did. A level-order array is scanned with
a queue; a depth-dashed preorder with a stack; a bracket string with recursion; two
traversal arrays need a split. ``Which reader?'' is decided entirely by the format, and
choosing the reader is the whole problem.

\begin{figure}[H]
\centering
\renewcommand{\arraystretch}{1.35}
\begin{tabular}{@{}p{0.30\linewidth} p{0.60\linewidth}@{}}
\toprule
\textbf{Encoding} & \textbf{The same tree $R$} \\
\midrule
Pointer tree & \code{TreeNode* root} (addresses; \S\ref{sec:c04-reps}) \\
Level-order array (LC) & \code{[1,2,3,4,5,\#,\#,\#,\#,6]} \\
Preorder $+$ null marks & \code{1,2,4,\#,\#,5,\#,\#,3,\#,6,\#,\#} \\
Nested brackets (LC 606) & \code{1(2(4)(5))(3()(6))} \\
Depth-dashed preorder & \code{1-2--4--5-3--6} \\
\bottomrule
\end{tabular}
\caption{One binary tree, five encodings. The \# in the level-order array marks \emph{skipped}
positions so later children line up; the \# in the preorder marks \emph{null children} so the
stream is self-delimiting (\thmref{c04-selfdelim}); $6$ sits under $3$'s \emph{right}
child, which the LC 606 brackets record explicitly as \code{(3()(6))} but the dashed preorder
cannot (see \probref{c04-recover-preorder}).}\label{fig:c04-encodings}
\end{figure}

Two design axes sort every encoding. First, \textbf{traversal order}: level-order (BFS)
arrays are flat and queue-friendly; preorder strings nest and recurse naturally. Second,
\textbf{how nulls are carried}: explicitly (preorder with \#), implicitly by position
(level-order array pads \# so children stay aligned), structurally (brackets \code{()}), or
not at all -- and \emph{``not at all'' is exactly when a single order stops being enough},
which is the content of \thmref{c04-selfdelim}.

\subsection{Definitions and Notation}

\begin{definition}{Encoding and decoding; self-delimiting}{c04-encoding}
An \term{encoding} of a tree is a string (or array) over a fixed alphabet from which the
tree can be \term{decoded} -- recovered -- by some reader. An encoding is
\term{self-delimiting} if, while scanning left to right, you always know when a subtree has
ended \emph{without} looking ahead; formally, no proper prefix of a valid encoding is itself
the encoding of a complete tree. Explicit null sentinels and balanced brackets make an
encoding self-delimiting.
\end{definition}

\noindent\textbf{In plain words.} A self-delimiting format tells you, token by token, exactly
where each node's children stop -- so one pass, no back-tracking, rebuilds the
tree. Formats that omit that signal (a bare preorder) leave you unable to tell where the left
subtree ends, hence not uniquely decodable.

\begin{definition}{LeetCode level-order serialisation}{c04-levelorder}
LeetCode writes a binary tree as the BFS queue of node values, \emph{including} a marker (\# /
\code{null}) for every missing child of a node that has been emitted, and trimming trailing
markers. Reading position $k$ (0-based) as a complete-tree slot, a node at slot $k$ has child
slots $2k+1$ and $2k+2$; the array is \emph{level} order, so a queue (not $2k$ arithmetic on a
skewed tree) is the safe reader.
\end{definition}

\subsection{Key Property: when is one order enough?}

\begin{theorem}{Unique decoding requires the nulls to be marked}{c04-selfdelim}
A preorder sequence of \emph{values alone} (no markers) determines a unique binary tree if and
only if null children are explicitly marked in the sequence, or a second order (typically the
inorder) is supplied. Without either, a bare preorder is consistent with \emph{many} trees.
\end{theorem}
\begin{proofbox}
($\Leftarrow$) With explicit null marks, decode recursively: read the root, then \emph{recurse
for the left subtree} and \emph{recurse for the right}; a mark stops a branch immediately. Each
subtree is bounded by its own marks, so the parse is forced -- the encoding is self-delimiting
and yields exactly one tree. ($\Rightarrow$) Conversely, given values-only preorder $1,2,3$: both
the tree ``$1$ with left child $2$ with left child $3$'' and ``$1$ with left $2$ and \emph{right}
$3$'' -- and three more -- produce the identical preorder $1,2,3$ (preorder lists node-then-
left-then-right, and cannot reveal which side a lone child occupies). Since the map
$\{\text{trees}\} \to \{\text{value-only preorders}\}$ is not injective, decoding is ambiguous
unless nulls are marked (adding information) or a second order pins each node's
side (inorder splits by the root -- \thmref{c03-root-position}, \secref{c03-rebuild}).
\end{proofbox}
\begin{insightbox}
This is precisely why LC 606 keeps an \emph{empty} \code{()} for a node with a right child but
no left (\probref{c04-string-tree}) -- dropping it would collapse \code{3()(6)} to \code{3(6)}
and lose the left/right distinction, exactly the ambiguity the proof constructs. The dashed
preorder of LC 1028 makes the \emph{opposite} choice (a lone child is always the left child) to
stay decodable with fewer tokens.
\end{insightbox}

\subsection{The Algorithm: two representative readers}

\begin{lstlisting}[caption={Decoding a level-order array, and emitting a general-tree bracket string}, label={lst:c04-formats}]
#include <bits/stdc++.h>
using namespace std;

struct TreeNode { int val; TreeNode *left, *right;
    TreeNode(int x=0): val(x), left(nullptr), right(nullptr) {} };

// READER 1 -- LeetCode level-order array ("1,2,3,#,4,5,#,#,#,#,#,6") -> TreeNode*
//   Queue-driven: pop a parent, attach the next two tokens as its children.
TreeNode* deserializeLevel(const vector<string>& tok) {
    if (tok.empty() || tok[0] == "#") return nullptr;
    TreeNode* root = new TreeNode(stoi(tok[0]));
    queue<TreeNode*> q; q.push(root);
    size_t i = 1;
    while (!q.empty() && i < tok.size()) {
        TreeNode* parent = q.front(); q.pop();
        if (i < tok.size() && tok[i] != "#") {              // left child
            parent->left = new TreeNode(stoi(tok[i])); q.push(parent->left);
        }
        ++i;
        if (i < tok.size() && tok[i] != "#") {              // right child
            parent->right = new TreeNode(stoi(tok[i])); q.push(parent->right);
        }
        ++i;
    }
    return root;
}

// READER 2 -- general tree -> nested bracket string "root(c1)(c2)..." (bracket encode)
//   DFS; wrap each child's rendering in parentheses (the LC-606 / Lisp shape).
string encodeBrackets(const vector<vector<int>>& kids, int u) {
    string s = to_string(u);
    for (int c : kids[u]) s += "(" + encodeBrackets(kids, c) + ")";
    return s;                                               // leaf -> just its label
}
\end{lstlisting}
\begin{complexitybox}
Both readers are $\bigO{n}$ in the number of tokens/nodes and $\bigO{n}$ space (queue /
recursion). Reader 1 uses a \emph{queue} -- the level-order array is a BFS object, so do not
try to read it as a preorder. Reader 2 recurses once per node; on a deep general tree this is
$\bigO{h}$ stack (chain risk as usual). The dash-preorder and the null-preorder readers
(\probref{c04-recover-preorder}) are $\bigO{n}$ with a stack instead.
\end{complexitybox}

\subsection{Worked Example}

\begin{example}{Decoding $R$'s level-order array}{c04-encodings}
Feed Reader 1 the array $[1,2,3,4,5,\#,\#,\#,\#,6]$ (\figref{c04-encodings}).
\textbf{Queue}: pop $1$, read tokens $2,3$ as its left,right (push both). Pop $2$, read $4,5$
(left,right, push). Pop $3$, read $\#$(no left) then $6$(right, push). Pop $4$, read $\#,\#$; pop
$5$, read nothing left. Result: $1\!\to\!\{2,3\}$, $2\!\to\!\{4,5\}$, $3\!\to\!\{\_,6\}$ -- exactly
$R$, with $6$ on $3$'s \emph{right}. The four $\#$ markers are what keep $4,5$ and
$6$ in their correct slots; drop one and everything shifts and the tree is wrong.
\end{example}
\begin{tracebox}[Reader 2 re-encodes $G$'s rooted children]
$\children$ of $G$ (from \secref{c04-rooting}): $0\!\to\!\{1,2,3\}$, $1\!\to\!\{4,5\}$, $3\!\to\!\{6\}$.
\code{encodeBrackets} on $0$: \code{"0"} + \code{"(" + enc(1) + ")"} + \code{"(" + enc(2) + ")"}
+ \code{"(" + enc(3) + ")"}. enc(1)$=$\code{1(4)(5)}, enc(2)$=$\code{2}, enc(3)$=$\code{3(6)}.
Result \code{"0(1(4)(5))(2)(3(6))}" -- the bracket encoding of \figref{c04-running-g}. \checkmark
\end{tracebox}

\subsection{Variations and Twists}

\begin{itemize}[nosep]
  \item \textbf{Multi-digit tokens}: a value may be several characters -- parse integers, do
        not index single chars (\probref{c04-recover-preorder} handles this).   \hfill\lv{3}
  \item \textbf{Comma vs space separators}: split on the delimiter before reading; a
        mis-set delimiter shifts every position.                               \hfill\lv{2}
  \item \textbf{JSON / nested-array input} (\code{[1,[2,[4]],[]]}\ldots): a bracket reader with
        a recursion + index, LC-flavoured.                                     \hfill\lv{4}
  \item \textbf{Lisp expression} (add/mul/let with scoping): a bracket parse \emph{plus} an
        environment stack.                                        \hfill\lv{5} (LC 736)
\end{itemize}

\subsection{Common Pitfalls}\label{sec:c04-formats-pitfalls}

\begin{warningbox}[Trap 1 -- making a node for a \# / null token]
In a level-order array the \# means ``this child slot is empty''. Creating a node for it gives
phantom zeros and a wrong shape; the check \code{if (tok[i] != "\#")} before \code{new TreeNode}
is mandatory.
\end{warningbox}
\begin{lstlisting}[caption={\textcolor{warningred}{WRONG}: allocates a node for the sentinel}, language=C++, basicstyle=\ttfamily\small]
parent->left = new TreeNode(stoi(tok[i]));   // stoi("#") is UB / throws
\end{lstlisting}
\begin{lstlisting}[caption={\textcolor{successgreen}{RIGHT}: only build when the token is a value}, language=C++, basicstyle=\ttfamily\small]
if (tok[i] != "#") { parent->left = new TreeNode(stoi(tok[i])); q.push(parent->left); }
\end{lstlisting}

\begin{warningbox}[Trap 2 -- reading a level-order array as if it were preorder]
The arrays carry different \emph{contracts}. Feeding $[1,2,3,\ldots]$ to a preorder decoder
attaches $3$ as $2$'s child instead of $1$'s right -- the tree is
mis-rooted. Match the reader to the format (\figref{c04-encodings}).
\end{warningbox}
\begin{lstlisting}[caption={\textcolor{warningred}{WRONG}: preorder recursion on a level-order array}, language=C++, basicstyle=\ttfamily\small]
// [1,2,3,4,...]: "root, then ALL of left subtree, then right" -- but this is BFS order!
\end{lstlisting}
\begin{lstlisting}[caption={\textcolor{successgreen}{RIGHT}: use the queue reader (Reader 1)}, language=C++, basicstyle=\ttfamily\small]
// pop parent, consume the NEXT TWO tokens as left,right -- level by level
\end{lstlisting}

\begin{warningbox}[Trap 3 -- dropping the empty \code{()} for a right-only child (LC 606)]
Encoding \code{3} with only a right child $6$ as \code{3(6)} instead of \code{3()(6)} loses the
left/right distinction and is \emph{not} uniquely decodable (\thmref{c04-selfdelim}). The empty
parentheses \emph{are} the information.
\end{warningbox}
\begin{lstlisting}[caption={\textcolor{warningred}{WRONG}: omits the left placeholder}, language=C++, basicstyle=\ttfamily\small]
s += "(" + tree2str(t->right) + ")";   // when left is null this mislabels right as left
\end{lstlisting}
\begin{lstlisting}[caption={\textcolor{successgreen}{RIGHT}: emit "()" for a null left when a right exists}, language=C++, basicstyle=\ttfamily\small]
s += "(" + tree2str(t->left) + ")";               // "()" if left is null
if (t->right) s += "(" + tree2str(t->right) + ")";
\end{lstlisting}

\subsection{Problem Lab}

One encoder and one decoder. LC 606 is a $\lv{2}$ \emph{output}-format problem whose only
teeth are the right-only-child \code{()} rule; LC 1028 is a genuine $\lv{4}$ decoder that
turns a depth-dashed preorder string back into pointers -- the inverse of a serialisation.


% ---------------------------------------------------------------------------
\solvedproblem{c04-string-tree}{Construct String from Binary Tree}{LeetCode 606}
  {https://leetcode.com/problems/construct-string-from-binary-tree/}{LC Easy L2}
  {preorder bracket encoding, the empty-parenthesis rule, encode (not decode)}

\begin{problemstatement}[Construct String from Binary Tree]
Given the root of a binary tree, build a string by \emph{preorder} traversal where every
node's value is printed, children are printed in parentheses, and \textbf{the empty
parentheses for a missing \emph{left} child must be printed only when a right child exists}.
Rules: a leaf prints just its number; a node with a left child prints
\code{val(leftsubtree)}; a node with \emph{only} a right child prints \code{val()(rightsubtree)}
(the \code{()} keeps the left slot). Example: tree $[1,2,3,4,\#,5,\#,6]$? no -- for the tree with
root $1$, left $2$ (left $4$), right $3$ (right $6$), the output is
\code{"1(2(4))(3()(6))"}. $1 \le n \le 3\cdot10^4$; values are non-negative integers.
\end{problemstatement}

\solstep{Understanding the Problem}
This is the \emph{forward} direction of \thmref{c04-selfdelim}: we are not parsing, we are
\textbf{emitting} a uniquely-readable bracket encoding of $R$-like trees. A preorder DFS that
wraps each child's rendering in \code{()} almost works -- the sole subtlety is the rule that
keeps left/right distinguishable: an empty \code{()} must appear for a missing left child
\emph{only if} there is a right child; a missing right child with a present left needs
nothing.

\solstep{Solution}
Recurse preorder. \code{nullptr} $\to$ \code{""}. A leaf $\to$ just its value. Otherwise print
\code{val}, then always \code{"(" + left + ")"} (which is \code{"()"} if left is \code{nullptr}),
then \code{"(" + right + ")"} \emph{only if} right exists.

\solstep{Implementation}
\begin{lstlisting}[caption={Construct String from Binary Tree (LeetCode 606)}, label={lst:c04-string-tree}]
#include <bits/stdc++.h>
using namespace std;

struct TreeNode { int val; TreeNode *left, *right;
    TreeNode(int x=0): val(x), left(nullptr), right(nullptr) {} };

class Solution {
public:
    string tree2str(TreeNode* t) {
        if (!t) return "";                                  // null contributes empty
        string s = to_string(t->val);
        if (!t->left && !t->right) return s;                // leaf: no parentheses
        s += "(" + tree2str(t->left) + ")";                 // left slot ALWAYS wrapped
        if (t->right) s += "(" + tree2str(t->right) + ")";  // right only when it exists
        return s;                                           // (null left + right exists -> "()")
    }
};
\end{lstlisting}

\solstep{Trace}
\begin{tracebox}[$R$ (node $3$ has no left child, right child $6$)]
Leaves: $4\to$\code{"4"}, $5\to$\code{"5"}, $6\to$\code{"6"}.
$2$ (left $4$, right $5$): \code{"2" + "(4)" + "(5)"} $=$ \code{"2(4)(5)"}.
$3$ (no left, right $6$): \code{"3" + "(())" $\to$ "()" + "(6)"} $=$ \code{"3()(6)"} -- the
empty \code{()} marks the absent left.
Root $1$: \code{"1" + "(2(4)(5))" + "(3()(6))"} $=$ \code{"1(2(4)(5))(3()(6))}". \checkmark\ \probref{c04-string-tree}.
\end{tracebox}

\solstep{Complexity}
\begin{complexitybox}
Time $\bigO{n}$ (each node visited once; total string length is $\bigO{n}$ including
parentheses), space $\bigO{h}$ recursion. Note \code{s += ...} building strings by value can be
$O(n^2)$ in the worst concatenation pattern if naively done at every level -- here each node
concatenates its own children once, so total is $\bigO{n}$. Verdict: \ac.
\end{complexitybox}

\begin{insightbox}[Key takeaway]
The empty \code{()} is not decoration -- it is the information that keeps the encoding
\emph{uniquely decodable} (\thmref{c04-selfdelim}). Encoding problems have all the same
formatting discipline as decoding ones, just read backwards.
\end{insightbox}

% ---------------------------------------------------------------------------
\solvedproblem{c04-recover-preorder}{Recover a Tree From Preorder Traversal}{LeetCode 1028}
  {https://leetcode.com/problems/recover-a-tree-from-preorder-traversal/}{LC Hard L4}
  {depth-dashed preorder, stack keyed by depth, multi-digit parse, decode}

\begin{problemstatement}[Recover a Tree From Preorder Traversal]
A binary tree is written in \emph{preorder}; before each node value there are exactly $d$
dashes \code{-} where $d$ is that node's depth (root has $0$). Values may be
\textbf{multi-digit}. Given the string \code{traversal}, rebuild the tree and return its root.
If a node has a single child, it is the \emph{left} child. Example: \code{"1-2--3--4-5--6--7"}
rebuilds root $1$, left $2$ (children $3,4$), right $5$ (children $6,7$). $1 \le$
length $\le 3\cdot10^4$; $1 \le$ node values $\le 10^9$.
\end{problemstatement}

\begin{attemptgate}[Attempt it yourself first -- \lv{4}, give yourself 30 minutes]
A fair first attempt: tokenise the string (dash-count = depth, then an integer), then think
about how a preorder-with-depth rebuilds.
\peek{when you read a node at depth $d$, its parent is the last node you saw at depth $d-1$;
what structure remembers ``the current node at each depth''?}
\peek{multi-digit: the number can be up to $10$ characters -- do not treat one char as one value.}
\nudge{keep a stack whose size equals the current path length; before placing a node, pop
until the stack holds exactly the ancestors (size $= d$); attach as parent's
left if free, else right.}
\fullhint{Single left-to-right pass: parse (depth, value) tokens; maintain a \code{stack} of
the current root-to-node path; pop to depth $d$, parent $=$ top, set left-else-right, push.
Root is the depth-$0$ node. Full build below.}
\end{attemptgate}

\solstep{Understanding the Problem}
The format is a \textbf{depth-annotated preorder}. Because preorder lists node-then-left-then-
right, the very next token is either the current node's left child (depth $+1$) or, after a
subtree finishes, a shallower node (depth $\le$ current). The dash count hands us the depth for
free, so the reader is: parse tokens, and let a stack track the current path -- when a node of
depth $d$ arrives, the $d$-th stack slot is its parent. This replaces the recursion of
\secref{c03-rebuild} with an explicit path stack and needs no inorder array.

\solstep{Key Observations}
\begin{observation}
A node at depth $d$ is attached to the most recent node at depth $d-1$ -- which is exactly the
top of a stack kept at size $d$ after popping deeper entries. Preorder guarantees that parent is
already built.
\end{observation}
\begin{observation}
Attach order encodes the side: the \emph{first} child seen under a parent is its left, the
\emph{second} is its right. A single child is therefore left (the problem's
convention) -- the format cannot otherwise tell left from right for a lone child.
\end{observation}
\begin{observation}
Values are multi-digit, so the tokenizer must greedily read all consecutive digits into one
integer (\code{val = val*10 + (c-'0')}), not emit one node per character.
\end{observation}

\solstep{Think It Through With Me}
\begin{thinkbox}
\deadend{Reparse dashes into an adjacency list and rebuild by rooting.}{There are no explicit
edges -- the dashes carry depth, and a tree is the \emph{preorder} of that rooting. Building an
edge list discards the visit order that tells you which node is whose parent; you would have to
re-derive the same stack logic anyway.}
\deadend{Branch over ``is this lone child left or right?''}{$2^{\text{\#lone-children}}$
possibilities -- exponential, and \emph{unnecessary}: LC fixes the convention (lone child $=$
left). The ambiguity the dash format has is \emph{resolved by the spec}, not searched over.}
\deadend{Use \code{split('-')} to get tokens.}{An empty token between consecutive dashes
encodes the depth, but \code{split} on \code{'-'} loses the count of \emph{how many} dashes
separated nodes -- depth is the number of dashes, so you must \emph{count} them, not split on
them.}
\aha{One pass with a path \emph{stack}: pop until size $= d$, parent $=$ top, attach
left-else-right, push the node -- depth handling is just stack height.}
\end{thinkbox}

\solstep{The Approach Ladder}
\begin{center}\renewcommand{\arraystretch}{1.3}
\begin{tabular}{@{}c p{0.235\linewidth} p{0.28\linewidth} l c@{}}
\toprule
\textbf{\#} & \textbf{Approach} & \textbf{Idea} & \textbf{Time} & \textbf{Verdict} \\
\midrule
1 & \lv{4} branch on lone-child side & try left/right for each single child & $O(2^{k})$ & \tle \\
2 & \lv{4} recursive descent by depth & consume a subtree given its depth & $O(n)$ & \ac \\
3 & \lv{4} path stack, single pass & pop to depth, attach left-else-right & $O(n)$ & \ac \\
\bottomrule
\end{tabular}\end{center}
Rungs 2 and 3 are both linear and correct; the stack (3) is written below because it avoids
$\bigO{h}$ recursion and makes the ``size $=$ depth'' invariant explicit.

\solstep{Solution}
Tokenise into (depth, value) pairs. Walk them with a \code{stack} of \code{TreeNode*}: pop
until the stack size equals the node's depth; make the node; if the stack is empty it is
the root, else attach it as the parent's left if that is free, else right; push the node.
Return the root.

\solstep{Proof of Correctness}
\begin{proofbox}[Loop invariant on the path stack]
\emph{Invariant:} after processing each token, the stack holds exactly the chain of nodes from
the root down to the most recently emitted node, i.e.\ the current root-to-frontier path, and
\code{stack[i]} has depth $i$. \emph{Base:} before any token the stack is empty; the first token
has depth $0$, popping to size $0$ is a no-op, its node becomes the root (empty stack), and is
pushed -- invariant holds. \emph{Step:} for a token at depth $d$, by the invariant all nodes at
depth $> d$ (a finished deeper subtree) are strictly above position $d$; popping to size $d$
discards exactly them, leaving \code{stack[d-1]} $=$ the last-seen depth-$(d-1)$ node $=$ the
parent (preorder guarantees the parent of a depth-$d$ node was the most recent depth-$(d-1)$
node). Left-else-right attachment reproduces preorder's ``left subtree before right'' and the
lone-child-left convention. Pushing the node restores the invariant (it is the new frontier).
By induction the built pointers form a tree whose preorder-with-depth is exactly the input
tokens, so the reconstruction is correct and unique under the convention.
\end{proofbox}

\solstep{Implementation}
\begin{lstlisting}[caption={Recover a Tree From Preorder Traversal (LeetCode 1028)}, label={lst:c04-recover-preorder}]
#include <bits/stdc++.h>
using namespace std;

struct TreeNode { int val; TreeNode *left, *right;
    TreeNode(int x=0): val(x), left(nullptr), right(nullptr) {} };

class Solution {
public:
    TreeNode* recoverFromPreorder(string traversal) {
        int n = (int)traversal.size(), i = 0;
        vector<pair<int,int>> tok;                   // (depth, value)
        while (i < n) {
            int d = 0;
            while (i < n && traversal[i] == '-') { ++d; ++i; }   // count dashes = depth
            long long val = 0;
            while (i < n && isdigit((unsigned char)traversal[i]))  // greedy multi-digit
                { val = val * 10 + (traversal[i] - '0'); ++i; }
            tok.push_back({d, (int)val});
        }

        stack<TreeNode*> path;                       // invariant: path[k] has depth k
        TreeNode* root = nullptr;
        for (auto [d, val] : tok) {
            while ((int)path.size() > d) path.pop(); // discard finished deeper subtree
            TreeNode* cur = new TreeNode(val);
            if (path.empty()) root = cur;            // depth 0 -> the root
            else {
                TreeNode* p = path.top();            // parent = last node at depth d-1
                if (!p->left) p->left = cur; else p->right = cur;   // preorder side rule
            }
            path.push(cur);
        }
        return root;
    }
};
\end{lstlisting}

\solstep{Trace}
\begin{tracebox}[$R$'s dashed preorder \code{1-2--4--5-3--6}, tokens $(0,1)(1,2)(2,4)(2,5)(1,3)(2,6)$]
$(0,1)$: path empty $\to$ root $1$, push $[1]$.
$(1,2)$: size $1 \ngtr 1$; parent $1$, left free $\to 1.\text{left}=2$, push $[1,2]$.
$(2,4)$: size $2 \ngtr 2$; parent $2$, left $\to 2.\text{left}=4$, push $[1,2,4]$.
$(2,5)$: pop to size $2$ (drop $4$) $[1,2]$; parent $2$, left used $\to 2.\text{right}=5$, push $[1,2,5]$.
$(1,3)$: pop to size $1$ (drop $5,2$) $[1]$; parent $1$, left used $\to 1.\text{right}=3$, push $[1,3]$.
$(2,6)$: size $2 \ngtr 2$; parent $3$, left free $\to 3.\text{left}=6$. \textbf{Note}: the format
places $6$ as $3$'s \emph{left} -- it cannot say ``right-only''; that is exactly
the ambiguity LC 606 resolves by emitting \code{()} (\thmref{c04-selfdelim}). The tree is $R$
mirrored at node $3$. \checkmark\ \probref{c04-recover-preorder}.
\end{tracebox}

\solstep{Complexity}
\begin{complexitybox}
Time $\bigO{L}$ where $L$ is the string length (each char scanned a constant number of times;
each node pushed and popped from the stack at most once). Space $\bigO{n}$ tokens plus
$\bigO{h}$ stack -- and the stack is \emph{heap} memory, so unlike a recursive descent it never
overflows the call stack on a $3\cdot10^4$-long chain. Verdict: \ac.
\end{complexitybox}

\solstep{Verification and Edge Cases}
\begin{center}\renewcommand{\arraystretch}{1.2}
\begin{tabular}{@{}p{0.30\linewidth} p{0.62\linewidth}@{}}
\toprule
\textbf{Case} & \textbf{What the code does} \\
\midrule
Single node \code{"10"} & one token $(0,10)$, root only, both children \code{nullptr} \\
Multi-digit like \code{1000000000} & greedy digit loop builds it into one value \\
Right-skewed chain & stack pops one each step; depth handling unchanged, no RE \\
Left-only children & every lone node attaches to \code{left} (the convention) \\
\bottomrule
\end{tabular}\end{center}

\solstep{Pitfalls and Variants}
\begin{warningbox}
Three bugs dominate: (i) reading one character as a value (multi-digit breaks); (ii)
\code{while(path.size() >= d)} instead of \code{> d} -- off-by-one that pops the parent; (iii)
attaching right before left (the preorder side order flips the tree). Variant: the identical
stack works for a \emph{general} tree (attach to the top's \code{children} list
instead of left/right) and for LC 938-style ``depth $+$ value'' encodings.
\end{warningbox}

\begin{insightbox}[Key takeaways -- three things to remember]
\begin{enumerate}[nosep, leftmargin=1.6em]
  \item A depth-annotated preorder decodes in one pass with a \textbf{path stack}: pop to the
        parent's depth, attach, push. Stack height \emph{is} the depth.
  \item Any string reader must respect \textbf{multi-digit} tokens -- count, then greedily parse
        digits.
  \item Encodings that omit the left/right signal for a lone child must fix a convention; LC 606
        keeps the signal (\code{()}), LC 1028 fixes it (lone $=$ left) -- \thmref{c04-selfdelim}.
\end{enumerate}
\end{insightbox}

\begin{hardeningbox}
\begin{itemize}[nosep, leftmargin=1.4em]
  \item \textbf{Recursion alternative}: a recursive-descent reader (consume subtree at depth
        $d$) is shorter but $\bigO{h}$ stack-deep -- on a skewed $3\cdot10^4$ tree the explicit
        \code{stack} above is the safe choice; keep it iterative.
  \item \textbf{Value overflow}: LC bounds values to $10^9$ (fits \code{int}), but read into
        \code{long long} while accumulating so a longer judge bound cannot silently wrap.
  \item \textbf{Memory}: emitting all $(d,\text{val})$ tokens up front is $O(n)$; you can also
        parse-and-place in a single fused pass to halve the constant when $L$ is large.
\end{itemize}
\end{hardeningbox}

\subsection{Practice Problems}

\begin{stuckbox}
\begin{itemize}[nosep, leftmargin=1.4em]
  \item Are you sure which \emph{order} the input array is -- level (BFS/queue) or preorder
        (DFS/stack)? (\figref{c04-encodings}, Trap 2.)
  \item In a string reader, does your tokenizer read a \emph{whole} integer, not one digit?
  \item When \emph{emitting} brackets, did you keep the empty \code{()} that preserves a
        right-only child? (Trap 3.)
  \item Does your depth/stack invariant pop \emph{to} the parent depth (\code{size $>$ d}), not
        too far?
\end{itemize}
\end{stuckbox}

\begin{practiceset}
  \practice{Score of Parentheses}{LeetCode 856}
    {https://leetcode.com/problems/score-of-parentheses/}{LC Medium L3}
    {A balanced-parentheses string \emph{is} a nested (bracket) encoding of a tree: depth $=$
    nesting level, each \code{()} a leaf. Track a running \code{score} with a depth counter (or a
    stack): on \code{(} go deeper (push $0$); on \code{)} pop, and add $\max(2\cdot\text{child},1)$
    to the new top. Decode-by-depth, one pass.}
  \practice{Parse Lisp Expression}{LeetCode 736}
    {https://leetcode.com/problems/parse-lisp-expression/}{LC Hard L4}
    {A Lisp-like nested encoding: \code{let}/\code{add}/\code{mul} with variable \emph{scoping}.
    Decode recursively, carrying a \textbf{stack of environments}; \code{let} pushes a scope of
    name/value bindings, \code{add}/\code{mul} recurse on operands, a bare var resolves in the
    innermost scope, and scopes pop as you return -- the same ``structure lives in nesting''
    reader as LC 1028, with an environment in each frame.}
\end{practiceset}

\begin{recallbox}[10-second refresher -- 4.3 Encodings and Input Formats]
Read the \emph{format}, then pick the reader: level-order array $\to$ \textbf{queue}; preorder
with \#/depth marks $\to$ \textbf{stack} or recursion; brackets $\to$ recursion. A single order is
uniquely decodable only when nulls are marked or a second order is given
(\thmref{c04-selfdelim}). When \emph{emitting}, the empty \code{()} is load-bearing.
\selftest{
  \item Write $R$ as (a) the level-order array, (b) the null-marked preorder, (c) the LC 606
        bracket string -- from memory.
  \item Why must the level-order reader use a queue and not a preorder recursion? (Trap 2.)
  \item State \thmref{c04-selfdelim} and give a value-only preorder that has two distinct trees.
  \item In LC 1028, why does a lone child become the \emph{left} child, and what would LC 606
        print differently?
  \item Parse the (depth, value) tokens of \code{1-2--3-4} and place them with the path stack.
}
\end{recallbox}
